\section{Complete theta source and finite evaluation}\label{sec:theta}

The de Bruijn--Newman field in physical height \(z\) is
\begin{equation}\label{eq:theta-source}
\begin{aligned}
H_\lambda(z)&=\int_0^\infty e^{\lambda u^2}\Phi(u)\cos(zu)\,du,\\
\Phi(u)&=\sum_{n=1}^\infty\Phi_n(u),\\
\Phi_n(u)&=\bigl(2\pi^2 n^4 e^{9u}-3\pi n^2 e^{5u}\bigr)
                     e^{-\pi n^2 e^{4u}}\qquad (u\ge 0).
\end{aligned}
\end{equation}
In this normalization \(H_0(z)=\xi(1/2+iz/2)/8\), with \(\xi\)
the completed Riemann xi function \cite[eqs.~(1)--(4)]{RodgersTao}.
The decay bound below makes \(H_\lambda\) entire in \(z\) and
justifies differentiation in \(\lambda\) on compact parameter sets.
Differentiating the heat weight and the cosine gives
\(\partial_\lambda H_\lambda=-\partial_{zz}H_\lambda\), hence
\(h_\tau=h_{zz}\) for \(h(\tau,z)=H_{\lambda_s-\tau}(z)\).

\subsection{Errors from the two omitted regions}

Take \(L=\lambda_s\), \(0\le \tau\le\Delta\le L\le 1/2\), an
integer \(N\ge 1\), and \(W\ge 1\). Use the finite model
\begin{equation}\label{eq:finite-theta-model}
P(\tau,x)=\sum_{n=1}^N \int_0^W
          e^{(L-\tau)u^2}\Phi_n(u)\cos(xu)\,du .
\end{equation}
Since
\(\Phi_n(u)=\pi n^2 e^{5u}(2\pi n^2 e^{4u}-3)e^{-\pi n^2 e^{4u}}>0\),
we can bound the absolute derivative error using the positive
source. For \(u\ge 0\) and \(n\ge M\ge 1\),
\[
\frac{(n+1)^4 e^{-\pi(n+1)^2 e^{4u}}}
     {n^4 e^{-\pi n^2 e^{4u}}}
\le r_M,\qquad
r_M=(1+1/M)^4 e^{-\pi(2M+1)}<1 .
\]
Indeed \(r_M\le r_1=16e^{-3\pi}\). In particular, summing from
\(M=1\) yields
\begin{equation}\label{eq:theta-density-majorant}
0<\Phi(u)\le C_\Phi e^{9u-\pi e^{4u}},
\qquad C_\Phi=\frac{2\pi^2}{1-16e^{-3\pi}} .
\end{equation}
The majorant dominates the source integrals with any fixed powers of
\(u\) and with \(e^{\operatorname{Re}(\lambda)u^2+|\operatorname{Im}z|u}\)
uniformly on compact sets. This gives absolute convergence and the
differentiability used in \eqref{eq:theta-source}.

For \(0\le k\le 6\), the \(k\)-th cosine derivative is bounded by
\(u^k\) for real \(x\). On \(0\le u\le W\) only the indices
\(n\ge M=N+1\) are omitted. Their positive source sum is at most
\[
\frac{2\pi^2M^4}{1-r_M}e^{9u-\pi M^2 e^{4u}}
\le \frac{2\pi^2M^4}{1-r_M}e^{-\pi M^2}.
\]
The last exponent \(9u-\pi M^2e^{4u}\) decreases from \(u=0\).
Integrating \(u^k\) and bounding the heat weight by \(e^{LW^2}\)
gives the omitted-index bound
\begin{equation}\label{eq:theta-index-tail}
A_k(N,W,L)=
\frac{2\pi^2M^4}{1-r_M}
e^{LW^2-\pi M^2}\frac{W^{k+1}}{k+1}.
\end{equation}

On \(u>W\) every index is omitted. Apply
\eqref{eq:theta-density-majorant} and put
\(\psi_k(u)=k\log u+Lu^2+9u-\pi e^{4u}\). Then
\[
\psi_k''(u)=-k/u^2+2L-16\pi e^{4u}<0,\qquad
-\psi_k'(W)=d_k(W,L):=4\pi e^{4W}-2LW-9-k/W>0 .
\]
For the last inequality, \(d_k(W,L)\ge 4\pi e^{4W}-W-15\),
which is positive at \(W=1\) and increasing thereafter.
The tangent inequality
\(\psi_k(u)\le\psi_k(W)-d_k(W,L)(u-W)\) bounds the entire
omitted integration region by
\begin{equation}\label{eq:theta-integration-tail}
T_k(W,L)=
\frac{C_\Phi W^k e^{LW^2+9W-\pi e^{4W}}}{d_k(W,L)} .
\end{equation}
The two regions are disjoint up to endpoints and exhaust the
error. Consequently
\begin{equation}\label{eq:theta-tail}
\bigl|\partial_x^k(h-P)(\tau,x)\bigr|\le
\epsilon_k:=A_k(N,W,L)+T_k(W,L)
\quad(0\le\tau\le\Delta,\ x\in\mathbb R,\ 0\le k\le 6).
\end{equation}
These real-axis allowances give the divided-difference errors
\(e_j=\epsilon_{j+2}/((j+1)(j+2))\) for \(0\le j\le 2\).
A complex argument would require a bound for the extra
\(e^{|\operatorname{Im}z|u}\) factor. Any numerical evaluation of
\eqref{eq:finite-theta-model} must also enclose quadrature and
rounding errors.

\subsection{Variation on parameter cells}

For \(0\le k\le 6\), an absolute moment allowance for both the
complete source and the finite model is
\begin{equation}\label{eq:theta-moment}
\int_0^\infty u^k e^{Lu^2}\Phi(u)\,du
\le \overline M_k(L):=\frac{C_\Phi e^{L-\pi}}{k+1}+T_k(1,L).
\end{equation}
On \([0,1]\) use \(e^{Lu^2}\le e^L\) and the decreasing
factor \(e^{9u-\pi e^{4u}}\le e^{-\pi}\) in
\eqref{eq:theta-density-majorant}. On \([1,\infty)\) use
\eqref{eq:theta-integration-tail} with \(W=1\). The model retains a
submeasure of the complete positive source, so the same allowance
bounds its derivatives.

In the simplex representation for \(U_j\) and \(V_j\), the argument
is \(\xi=c+d\beta\), where
\[
\beta=-(1-s-t)+t+sv=-1+s+2t+sv,\qquad |\beta|\le 1 .
\]
The last inequality follows because \(\beta\) is a convex combination
of \(-1,1,v\). For \(Y_j=U_j\) or \(V_j\) and \(j=0,1,2\), the
variation on the box \eqref{eq:box} satisfies
\begin{equation}\label{eq:theta-cell}
\begin{aligned}
|\partial_c Y_j|,\ |\partial_d Y_j|
    &\le \frac{\overline M_{j+3}(L)}{(j+1)(j+2)},\\
|\partial_v Y_j|
    &\le \frac{D\,\overline M_{j+3}(L)}{(j+2)(j+3)},\\
|\partial_\tau Y_j|
    &\le \frac{\overline M_{j+4}(L)}{(j+1)(j+2)} .
\end{aligned}
\end{equation}
The \(c,d\) derivatives cost one spatial derivative and multiply
the integrand by \(1,\beta\), respectively. The \(v\) derivative
adds the factor \(ds\) to the simplex weight \(s^j\), replacing its
mass by at most \(D/((j+2)(j+3))\). At fixed \(c,d,v\) the
\(\tau\) derivative inserts \(-u^2\) into the transform, equivalent
to two spatial derivatives. Thus a cell allowance is the sum of
these derivative bounds times the corresponding coordinate radii.
The moment and cell bounds control variation and truncation. A
certified positive factor floor and an upper bound for the signed
curvature average require enclosures of the model on the full box.

\subsection{What the alternative primitive gives}

At a putative double event define \(C_j,S_j\) to be the cosine and
sine moments of \(u^j e^{\lambda_*u^2}\Phi(u)\) at \(x_*\).
Differentiating \eqref{eq:theta-source} gives
\(C_0=S_1=0\) and
\((h_{xx},h_{xxx},h_{xxxx})=(-C_2,S_3,C_4)\). The even moments
retain the oscillatory cosine. Absolute moment bounds therefore
leave the event's denominator and its signed curvature unpaid.

The inherited half-line primitive provides a second exact source
interface \cite{OARecord}. Put
\[
 p(u)=e^u\sum_{n\ge1}e^{-\pi n^2e^{4u}},\qquad
 J_\lambda(x)=\int_0^\infty p(u)e^{\lambda u^2}\cos(xu)\,\dd u.
\]
For \(y=\pi n^2e^{4u}\) each term \(e^u e^{-y}\) satisfies
\((\partial_u^2-1)e^u e^{-y}=(16y^2-24y)e^u e^{-y}\).
Hence \(p''-p=8\Phi\). Theta reflection gives \(p'(0)=-1/2\):
if \(\theta(w)=\sum_{n\in\mathbb Z}e^{-\pi n^2w}\), then
\(e^u\theta(e^{4u})\) is even by Poisson summation, and
\(2p(u)=e^u\theta(e^{4u})-e^u\). Two integrations by parts
give \(8\int_0^\infty\Phi f=f(0)/2+p(0)f'(0)+\int_0^\infty p(f''-f)\)
for the cosine test and its derivatives. The boundary terms at
infinity vanish by the same decay. Substituting
\(f=e^{\lambda u^2}\cos(xu)\) yields
\begin{equation}
8H_\lambda=\tfrac12+(2\lambda-1-x^2)J_\lambda+
                  4\lambda xJ_\lambda'-4\lambda^2J_\lambda''.
\label{eq:primitive}
\end{equation}
Differentiating \(n=2,3,4\) times in \(x\) gives
\[
\begin{aligned}
8H_\lambda^{(n)}
={}&-4\lambda^2J_\lambda^{(n+2)}
    +4\lambda xJ_\lambda^{(n+1)}
    +(2\lambda-1-x^2+4\lambda n)J_\lambda^{(n)}\\
 &-2nxJ_\lambda^{(n-1)}-n(n-1)J_\lambda^{(n-2)} .
\end{aligned}
\]
Thus the primitive supplies the event jet through derivatives of
order six. Its positivity alone supplies neither a lower
denominator margin nor an interval estimate for \(B\).
